\sectionheading{5.18 প্রোগ্রাম ডিজাইনের কৌশল (Program Design Techniques)}

কোনো Problem-এর Solution সরাসরি Code করা শুরু করলে বড় Program
দ্রুত জটিল হয়ে যেতে পারে।

তাই Program তৈরির আগে সমস্যাটি কীভাবে বিভিন্ন অংশে সাজানো হবে
এবং কোন অংশ আগে তৈরি করা হবে, তা পরিকল্পনা করা প্রয়োজন।

এই পরিকল্পিত পদ্ধতিগুলোকে
\textbf{Program Design Techniques} বলা হয়।

%================================================
% Definition
%================================================
\booksubheading{Program Design Technique কাকে বলে?}

\begin{definitionbox}

\textbf{\color{darkTeal}সংজ্ঞা}

কোনো Problem-এর Solution-কে সুশৃঙ্খলভাবে পরিকল্পনা করে
Program-এর বিভিন্ন অংশ কীভাবে গঠন ও বাস্তবায়ন করা হবে,
তার পদ্ধতিকে \textbf{Program Design Technique} বলে।

\end{definitionbox}

\textbf{সহজভাবে মনে রাখবে—}

\begin{center}

\textbf{
Problem
\(\longrightarrow\)
পরিকল্পনা
\(\longrightarrow\)
অংশে বিন্যাস
\(\longrightarrow\)
Program তৈরি
}

\end{center}

Program Design-এর দুটি গুরুত্বপূর্ণ কৌশল হলো—

\begin{itemize}

\item \textbf{Top-down Approach}

\item \textbf{Bottom-up Approach}

\end{itemize}

%================================================
% Top-down
%================================================
\booksubheading{১. Top-down Approach}

\begin{definitionbox}

যে পদ্ধতিতে প্রথমে \textbf{সম্পূর্ণ Problem বা Main Task}
বিবেচনা করা হয় এবং পরে সেটিকে ধাপে ধাপে ছোট ছোট
Subproblem বা Module-এ ভাগ করা হয়, তাকে
\textbf{Top-down Approach} বলে।

\end{definitionbox}

\textbf{সহজভাবে—}

\begin{center}
\begin{adjustbox}{max width=0.98\linewidth}
\fcolorbox{darkBlue}{softBlue}{
    \parbox[c][9mm][c]{0.22\linewidth}{
        \centering\textbf{বড়\\Problem}
    }
}
\quad
\(\longrightarrow\)
\quad
\fcolorbox{mainTeal}{softGreen}{
    \parbox[c][9mm][c]{0.23\linewidth}{
        \centering\textbf{ছোট ছোট\\Problem}
    }
}
\quad
\(\longrightarrow\)
\quad
\fcolorbox{darkBlue}{softBlue}{
    \parbox[c][9mm][c]{0.22\linewidth}{
        \centering\textbf{আরও ছোট\\Task}
    }
}
\end{adjustbox}
\end{center}

অর্থাৎ Main Problem থেকে কাজ শুরু করে প্রয়োজন অনুযায়ী
প্রতিটি অংশকে আরও ছোট অংশে ভাগ করা হয়।

%================================================
% Stepwise Refinement
%================================================
\booksubheading{Stepwise Refinement}

Top-down Design-এ একটি বড় কাজকে ধাপে ধাপে আরও ছোট ও
বিস্তারিত কাজে ভাঙার প্রক্রিয়াকে
\textbf{Stepwise Refinement} বলা হয়।

\begin{center}

\textbf{
Main Problem
\(\longrightarrow\)
Subproblem
\(\longrightarrow\)
আরও ছোট Task
}

\end{center}

\begin{rememberbox}

\textbf{\color{deepOrange}MCQ-এর জন্য মনে রাখবে}

\textbf{Stepwise Refinement} মূলত
\textbf{Top-down Approach}-এর সঙ্গে সম্পর্কিত।

\end{rememberbox}

%================================================
% Top-down Example
%================================================
\begin{examplebox}{Top-down Approach-এর উদাহরণ}

একটি \textbf{Student Result Program} তৈরি করতে হলে প্রথমে
পুরো Problem-টিকে এভাবে ভাগ করা যেতে পারে—

\begin{center}

\textbf{
Student Result Program
}

\(\downarrow\)

\textbf{
Input Marks
\quad + \quad
Calculate Result
\quad + \quad
Display Result
}

\end{center}

আবার \textbf{Calculate Result} অংশকে আরও ভাগ করা যায়—

\begin{center}

\textbf{
Total Calculate
\quad + \quad
Average Calculate
\quad + \quad
Grade Determine
}

\end{center}

অর্থাৎ বড় Problem ধাপে ধাপে ছোট ও সহজ কাজে বিভক্ত হয়েছে।

\end{examplebox}

%================================================
% Bottom-up
%================================================
\booksubheading{২. Bottom-up Approach}

\begin{definitionbox}

যে পদ্ধতিতে প্রথমে প্রয়োজনীয়
\textbf{ছোট ছোট Module বা Component} তৈরি করা হয় এবং পরে
সেগুলোকে একত্রিত করে সম্পূর্ণ Program তৈরি করা হয়,
তাকে \textbf{Bottom-up Approach} বলে।

\end{definitionbox}

\textbf{সহজভাবে—}

\begin{center}
\begin{adjustbox}{max width=0.98\linewidth}
\fcolorbox{darkBlue}{softBlue}{
    \parbox[c][9mm][c]{0.23\linewidth}{
        \centering\textbf{ছোট ছোট\\Module}
    }
}
\quad
\(\longrightarrow\)
\quad
\fcolorbox{mainTeal}{softGreen}{
    \parbox[c][9mm][c]{0.22\linewidth}{
        \centering\textbf{Module\\Combine}
    }
}
\quad
\(\longrightarrow\)
\quad
\fcolorbox{darkBlue}{softBlue}{
    \parbox[c][9mm][c]{0.22\linewidth}{
        \centering\textbf{সম্পূর্ণ\\Program}
    }
}
\end{adjustbox}
\end{center}

%================================================
% Bottom-up Example
%================================================
\begin{examplebox}{Bottom-up Approach-এর উদাহরণ}

Student Result Program-এর ক্ষেত্রে প্রথমে আলাদাভাবে—

\begin{itemize}

\item Input Module
\item Total Calculation Module
\item Grade Module
\item Output Module

\end{itemize}

তৈরি ও পরীক্ষা করা হলো।

এরপর Moduleগুলো একত্রিত করে সম্পূর্ণ
\textbf{Student Result Program} তৈরি করা হলো।

\end{examplebox}

%================================================
% Comparison
%================================================
\booksubheading{Top-down ও Bottom-up Approach-এর তুলনা}

\begin{tabularx}{\linewidth}{|>{\bfseries}p{0.22\linewidth}|Y|Y|}
\hline

\rowcolor{softBlue}
\multicolumn{1}{|c|}{\textbf{বিষয়}}
&
\multicolumn{1}{c|}{\textbf{Top-down Approach}}
&
\multicolumn{1}{c|}{\textbf{Bottom-up Approach}}
\\
\hline

শুরু হয়
&
সম্পূর্ণ Problem থেকে
&
ছোট Module বা Component থেকে
\\
\hline

অগ্রগতি
&
বড় থেকে ছোট
&
ছোট থেকে বড়
\\
\hline

মূল কাজ
&
Problem-কে ছোট অংশে ভাগ করা
&
ছোট Moduleগুলো একত্র করা
\\
\hline

গুরুত্ব
&
Overall Design ও Decomposition
&
Module তৈরি ও Integration
\\
\hline

সহজ সূত্র
&
Whole \(\rightarrow\) Parts
&
Parts \(\rightarrow\) Whole
\\
\hline

\end{tabularx}

%================================================
% Relationship with Modular Programming
%================================================
\booksubheading{Modular Programming-এর সঙ্গে সম্পর্ক}

Program Design Technique এবং Modular Programming
পরস্পরের সঙ্গে সম্পর্কিত হলেও একই বিষয় নয়।

\textbf{Top-down Approach}-এ বড় Problem-কে ছোট Module-এ
ভাগ করা যায়।

\textbf{Bottom-up Approach}-এ ছোট Moduleগুলোকে যুক্ত করে
বড় Program তৈরি করা যায়।

\begin{center}

\textbf{
Top-down
\(\longrightarrow\)
Divide
}

\vspace{0.5mm}

\textbf{
Bottom-up
\(\longrightarrow\)
Combine
}

\end{center}

আগের Topic-এ Modular Programming বিস্তারিত আলোচনা করা হয়েছে,
তাই এখানে সেটি পুনরায় বিস্তারিত করা হলো না।

%================================================
% When Useful
%================================================
\booksubheading{কোন পদ্ধতি কখন সুবিধাজনক?}

\textbf{\color{darkTeal}Top-down Approach সুবিধাজনক যখন—}

\begin{itemize}

\item প্রথমে পুরো Problem-এর Structure বোঝা প্রয়োজন।

\item বড় Problem-কে ছোট অংশে ভাগ করতে হয়।

\item Main Task থেকে বিস্তারিত Design তৈরি করতে হয়।

\end{itemize}

\textbf{\color{darkTeal}Bottom-up Approach সুবিধাজনক যখন—}

\begin{itemize}

\item আগে থেকেই কিছু Module বা Component রয়েছে।

\item ছোট অংশগুলো আলাদাভাবে তৈরি ও Test করা সুবিধাজনক।

\item Reusable Component ব্যবহার করে বড় Program তৈরি করা হয়।

\end{itemize}

বাস্তব Program Development-এ প্রয়োজন অনুযায়ী
দুটি পদ্ধতির ধারণাই ব্যবহার করা যেতে পারে।

%================================================
% Quick Comparison
%================================================
\booksubheading{দ্রুত মনে রাখার তুলনা}

\begin{center}

\fcolorbox{darkBlue}{softBlue}{
    \parbox[c][13mm][c]{0.40\linewidth}{
        \centering
        \textbf{Top-down}\\[0.6mm]
        Whole \(\rightarrow\) Divide \(\rightarrow\) Submodules
    }
}
\qquad
\fcolorbox{mainTeal}{softGreen}{
    \parbox[c][13mm][c]{0.40\linewidth}{
        \centering
        \textbf{Bottom-up}\\[0.6mm]
        Modules \(\rightarrow\) Combine \(\rightarrow\) Whole
    }
}

\end{center}

%================================================
% Final Reminder
%================================================
\begin{rememberbox}

\textbf{\color{deepOrange}মনে রাখবে}

\begin{itemize}

\item \textbf{Program Design Technique} হলো Program তৈরির
পরিকল্পিত পদ্ধতি।

\item প্রধান দুটি কৌশল—
\textbf{Top-down এবং Bottom-up Approach}।

\item \textbf{Top-down = বড় থেকে ছোট}।

\item \textbf{Bottom-up = ছোট থেকে বড়}।

\item Top-down-এ Problem-কে ছোট অংশে
\textbf{ভাগ করা হয়}।

\item Bottom-up-এ ছোট Module বা Component
\textbf{একত্র করা হয়}।

\item \textbf{Stepwise Refinement} Top-down Approach-এর
সঙ্গে সম্পর্কিত।

\item Modular Programming-এর সঙ্গে দুই Design Technique-এর
ঘনিষ্ঠ সম্পর্ক রয়েছে।

\item সবচেয়ে সহজ Shortcut—

\begin{center}

\textbf{
Top-down = Divide
\qquad\qquad
Bottom-up = Combine
}

\end{center}

\end{itemize}

\end{rememberbox}

\small

%================================================
% 5.19 Computational Thinking
%================================================
